\subsection{Testing the Model on SPARC Data}

To verify the performance of Baseline v2 with the fixed theoretical parameter \( P \equiv \sqrt{8/3} \approx 1.63299 \) and the strict implementation of the decoupling force \(\mathbf{F}_{\rm decouple}\), a simplified test was conducted on the full sample of 175 galaxies from the SPARC catalog.

The test used characteristic parameters of each galaxy (total baryonic mass and maximum radius). For each galaxy, a patch simulation was performed using the new strict form of the decoupling projection operator.

\subsubsection{Test Results}

\begin{itemize}
    \item Galaxies processed: 175
    \item LSB galaxies: 61 (34.9\%)
    \item Average patch stability: \( 97.68\% \pm 0.42\% \)
    \item Median stability: 97.71\%
    \item Minimum stability: 96.12\%
    \item Average effective gravitational factor: 1.6447
    \item Theoretical parameter \( P \): 1.63299
\end{itemize}

Breakdown by galaxy type:
\begin{itemize}
    \item Normal galaxies (114): average stability \( 97.89\% \)
    \item LSB galaxies (61): average stability \( 97.28\% \)
    \item Difference between normal and LSB galaxies: \( \approx 0.61\% \)
\end{itemize}

The obtained results show that the transition to the strictly fixed value \( P = \sqrt{8/3} \) did not lead to a significant loss in the model's descriptive power. The model maintains high patch stability (above 96\% even for low-surface-brightness galaxies) while completely eliminating the fitting of the parameter \( P \).

The relatively small difference in stability between normal and LSB galaxies is explained by the simplified nature of the test, which used only integral galaxy parameters. A more pronounced difference is expected when using full radial profiles from the Rotmod\_LTG files.

These results confirm the validity of the chosen theoretical value of \( P \) and the functionality of the strict decoupling operator within Baseline v2.
